% Adapted from academic-drawing's TikZ template; project fonts and palette prevail.
% Final canvas: 167 mm wide. Ports: text east -> spans west; spans east
% -> cluster west; cluster east -> KB west. Side branches use a routing lane.
\begin{tikzpicture}[
  x=1mm,y=1mm,
  font=\figurefont\fontsize{9}{11}\selectfont,
  box/.style={rectangle,align=left,inner sep=2mm,line width=0.85pt,
    draw=BookRule,fill=white},
  source/.style={box,text width=31mm,draw=FigureInput,fill=FigureInput!8},
  span/.style={box,text width=27mm,draw=FigureInput,fill=FigureInput!8},
  cluster/.style={box,text width=23mm,draw=FigureModel,fill=FigureModel!8},
  target/.style={box,text width=33mm,draw=FigureOutput,fill=FigureOutput!8},
  locate/.style={draw=FigureInput,line width=0.85pt,rounded corners=1.2mm},
  coref/.style={draw=FigureModel,line width=0.85pt,dashed,rounded corners=1.2mm},
  link/.style={draw=FigureOutput,line width=1.1pt,-{Latex[length=2mm]}},
  head/.style={font=\figurefont\bfseries\fontsize{9}{11}\selectfont,
    text=BookInk,align=center}
]
  \node[head] at (18,37) {补充原文};
  \node[head] at (62,37) {原文跨度};
  \node[head] at (104,37) {文档共指簇};
  \node[head] at (146,37) {教学知识库};
  \node[source] (text) at (18,0) {
    地点为海棠楼201。\\
    该教室旁设临时服务点。\\[3pt]
    字符从0开始计数
  };
  \node[span] (m1) at (62,23) {海棠楼201\\$[3,9)$};
  \node[span] (m2) at (62,0) {该教室\\$[10,13)$};
  \node[span] (m3) at (62,-23) {临时服务点\\$[15,20)$};
  \node[cluster] (c1) at (104,11.5) {同一教室\\$\{m_1,m_2\}$};
  \node[cluster] (c2) at (104,-23) {服务点\\$\{m_3\}$};
  \node[target] (kb1) at (146,11.5) {ROOM-201\\海棠楼201};
  \node[target] (kb2) at (146,-23) {NIL\\无可确认库中项};
  \draw[locate] (text.east) -- (42,0);
  \draw[locate] (42,23) -- (42,-23);
  \draw[locate] (42,23) -- (m1.west);
  \draw[locate] (42,0) -- (m2.west);
  \draw[locate] (42,-23) -- (m3.west);
  \draw[coref] (m1.east) -- (84,23) -- (84,11.5) -- (c1.west);
  \draw[coref] (m2.east) -- (84,0) -- (84,11.5);
  \draw[coref] (m3.east) -- (c2.west);
  \draw[link] (c1.east) -- (kb1.west);
  \draw[link] (c2.east) -- (kb2.west);
  \draw[locate] (4,-38) -- (14,-38);
  \node[anchor=west] at (16,-38) {定位};
  \draw[coref] (49,-38) -- (59,-38);
  \node[anchor=west] at (61,-38) {共指归组};
  \draw[link] (111,-38) -- (121,-38);
  \node[anchor=west] at (123,-38) {身份选择};
\end{tikzpicture}
